% source: RKF/theorum/23_native_cut_generated_object_theorem.md
\hypertarget{native-cut-generated-object-theorem}{%
\chapter{Native Cut-Generated Object Theorem}\label{native-cut-generated-object-theorem}}

\hypertarget{purpose}{%
\section{1. Purpose}\label{purpose}}

This theorem makes the cut, not a preassembled operator, the primitive object.

The classical world may suggest that the final maps should resemble a source
operator, a boundary row, a Birman--Schwinger matrix, a Fredholm determinant or
a spectral-flow index. None of those objects is imported here. They are
constructed as compressions of one native event lift by the recognized and cut
projectors.

The theorem has two layers:

\begin{verbatim}
full cut covariance          possibly infinite dimensional;
target-relevant cut memory   dimension at most five after faithfulness.
\end{verbatim}

Confusing those layers would falsely claim that an infinite-dimensional
positive covariance has rank five. Humanity already has enough false miracles.

\begin{center}\rule{0.5\linewidth}{0.5pt}\end{center}

\hypertarget{primitive-data}{%
\section{2. Primitive data}\label{primitive-data}}

Let \(\mathcal F_0\) be the native generator core and \(\mathcal K\) the event
carrier. Let

\[
Z:\mathcal F_0\longrightarrow\mathcal K
\]

be the native event lift constructed from the declared cut events.

Let

\[
P=P^*=P^2,
\qquad
Q=I-P,
\qquad
PQ=QP=0
\]

be the recognized and cut-memory projectors.

Define the cut involution

\[
\mathfrak j=P-Q.
\]

Then

\[
\mathfrak j^*=\mathfrak j,
\qquad
\mathfrak j^2=I.
\]

\begin{center}\rule{0.5\linewidth}{0.5pt}\end{center}

\hypertarget{derived-objects}{%
\section{3. Derived objects}\label{derived-objects}}

Define

\[
S=Z^*Z,
\qquad
R=Z^*PZ,
\qquad
D=Z^*QZ,
\qquad
F=Z^*\mathfrak jZ.
\]

\hypertarget{theorem-3.1-cut-generated-decomposition}{%
\subsection{Theorem 3.1 (Cut-generated decomposition)}\label{theorem-3.1-cut-generated-decomposition}}

The four forms satisfy

\[
\boxed{S=R+D}
\]

and

\[
\boxed{F=R-D.}
\]

Moreover,

\[
R\ge0,
\qquad
D\ge0.
\]

\hypertarget{proof}{%
\subsection{Proof}\label{proof}}

Since \(P+Q=I\),

\[
S=Z^*(P+Q)Z=Z^*PZ+Z^*QZ=R+D.
\]

Since \(\mathfrak j=P-Q\),

\[
F=Z^*(P-Q)Z=R-D.
\]

The positivity of \(R\) and \(D\) follows because they are Gram forms:

\[
\langle f,Rf\rangle=\|PZf\|^2,
\qquad
\langle f,Df\rangle=\|QZf\|^2.
\]

∎

Thus source energy, recognized energy, cut-memory energy and signed cut form are
co-generated by one cut lift.

\begin{center}\rule{0.5\linewidth}{0.5pt}\end{center}

\hypertarget{first-moment-and-second-moment}{%
\section{4. First moment and second moment}\label{first-moment-and-second-moment}}

Assume that the recognized projector has rank one,

\[
P=|p\rangle\langle p|,
\qquad
\|p\|=1.
\]

Define the native first moment

\[
L(f)=\langle p,Zf\rangle.
\]

\hypertarget{theorem-4.1-native-cut-covariance}{%
\subsection{Theorem 4.1 (Native cut covariance)}\label{theorem-4.1-native-cut-covariance}}

The recognized Gram is

\[
R=L^*L,
\]

and therefore

\[
\boxed{S-L^*L=D=Z^*QZ\ge0.}
\]

For \(f,h\in\mathcal F_0\),

\[
\boxed{
S(f,h)-\overline{L(h)}L(f)
=
\langle QZh,QZf\rangle.
}
\]

\hypertarget{proof-1}{%
\subsection{Proof}\label{proof-1}}

Since \(P=|p\rangle\langle p|\),

\[
R(f,h)
=
\langle Zh,PZf\rangle
=
\langle Zh,p\rangle\langle p,Zf\rangle
=
\overline{L(h)}L(f).
\]

Apply Theorem 3.1. ∎

This is the operator form of the Cut-Square Identity. Finite coordinate
identities are shadows of this projection formula.

\begin{center}\rule{0.5\linewidth}{0.5pt}\end{center}

\hypertarget{exact-exponential-transport}{%
\section{5. Exact exponential transport}\label{exact-exponential-transport}}

Let

\[
U_t=e^{itG}
\]

be the exact finite-time exponential change generated by a self-adjoint native
change generator \(G\). Transport the event lift and cut by

\[
Z_t=U_tZ_0,
\qquad
P_t=U_tP_0U_t^*,
\qquad
Q_t=U_tQ_0U_t^*.
\]

\hypertarget{theorem-5.1-exponential-cut-covariance}{%
\subsection{Theorem 5.1 (Exponential cut covariance)}\label{theorem-5.1-exponential-cut-covariance}}

For every \(t\),

\[
Z_t^*P_tZ_t=Z_0^*P_0Z_0,
\]

\[
Z_t^*Q_tZ_t=Z_0^*Q_0Z_0,
\]

and

\[
Z_t^*(P_t-Q_t)Z_t=Z_0^*(P_0-Q_0)Z_0.
\]

\hypertarget{proof-2}{%
\subsection{Proof}\label{proof-2}}

Use \(U_t^*U_t=I\). For example,

\[
Z_t^*Q_tZ_t
=Z_0^*U_t^*U_tQ_0U_t^*U_tZ_0
=Z_0^*Q_0Z_0.
\]

The other identities are identical. ∎

Thus analyticity is not imposed by a classical continuation theorem. It is
exact covariance of the cut-generated objects under finite exponential change.

\begin{center}\rule{0.5\linewidth}{0.5pt}\end{center}

\hypertarget{target-relevant-memory}{%
\section{6. Target-relevant memory}\label{target-relevant-memory}}

Let

\[
C=QZ
\]

be the full cut bridge. It may have infinite-dimensional range.

Let

\[
T_\Sigma:\mathcal K\longrightarrow\mathbb C^\nu,
\qquad
\nu\le5,
\]

be the native target observer, with

\[
T_\Sigma T_\Sigma^*=I_\nu.
\]

Define

\[
A_\Sigma=T_\Sigma C.
\]

\hypertarget{definition-6.1-target-faithfulness}{%
\subsection{Definition 6.1 (Target faithfulness)}\label{definition-6.1-target-faithfulness}}

The target observer is faithful on the adverse cut module if

\[
\boxed{C=T_\Sigma^*T_\Sigma C.}
\]

The hidden residual is

\[
H_\Sigma=(I-T_\Sigma^*T_\Sigma)C.
\]

Thus target faithfulness is exactly \(H_\Sigma=0\) on the module being
reduced.

\hypertarget{theorem-6.2-five-memory-reduction}{%
\subsection{Theorem 6.2 (Five-memory reduction)}\label{theorem-6.2-five-memory-reduction}}

Under target faithfulness,

\[
\operatorname{rank}C\le\nu\le5
\]

on the target-relevant adverse module, and

\[
\boxed{D=C^*C=A_\Sigma^*A_\Sigma.}
\]

The nonzero spectra of \(D\) and

\[
B_\Sigma^{\rm cut}=A_\Sigma A_\Sigma^*
\]

coincide with multiplicity.

\hypertarget{proof-3}{%
\subsection{Proof}\label{proof-3}}

Target faithfulness gives

\[
C=T_\Sigma^*A_\Sigma.
\]

Hence the range of \(C\) factors through \(\mathbb C^\nu\). Also,

\[
C^*C
=A_\Sigma^*T_\Sigma T_\Sigma^*A_\Sigma
=A_\Sigma^*A_\Sigma.
\]

The spectral statement follows from the cut-bridge spectral isomorphism. ∎

\begin{center}\rule{0.5\linewidth}{0.5pt}\end{center}

\hypertarget{hidden-complement-theorem}{%
\section{7. Hidden complement theorem}\label{hidden-complement-theorem}}

A target observer need not see all cut memory. The exact decomposition is

\[
C=T_\Sigma^*A_\Sigma+H_\Sigma,
\]

with orthogonal ranges, and hence

\[
D=A_\Sigma^*A_\Sigma+H_\Sigma^*H_\Sigma.
\]

\hypertarget{theorem-7.1-no-hidden-adverse-complement-criterion}{%
\subsection{Theorem 7.1 (No-hidden-adverse-complement criterion)}\label{theorem-7.1-no-hidden-adverse-complement-criterion}}

Suppose the signed endpoint form has the cut-derived form

\[
\mathcal W
=R_0-A_\Sigma^*A_\Sigma+H_\Sigma^*H_\Sigma,
\qquad
R_0>0.
\]

Then the hidden complement is favorable. Therefore

\[
R_0-A_\Sigma^*A_\Sigma\ge0
\quad\Longrightarrow\quad
\mathcal W\ge0.
\]

If \(H_\Sigma=0\), the finite reduction is exact.

\hypertarget{proof-4}{%
\subsection{Proof}\label{proof-4}}

The term \(H_\Sigma^*H_\Sigma\) is positive. ∎

This theorem identifies the precise burden. The framework need not prove that
all cut memory is five-dimensional. It must prove only that any target-hidden
complement is nonnegative.

\begin{center}\rule{0.5\linewidth}{0.5pt}\end{center}

\hypertarget{native-finite-seam-matrix}{%
\section{8. Native finite seam matrix}\label{native-finite-seam-matrix}}

Define

\[
V_\Sigma=A_\Sigma^*:
\mathbb C^\nu\longrightarrow\mathcal F^{\rm comp}
\]

and, for a positive cut-derived reference \(R_0\),

\[
B_\Sigma
=V_\Sigma^*R_0^{-1}V_\Sigma
=A_\Sigma R_0^{-1}A_\Sigma^*.
\]

Then

\[
R_0-A_\Sigma^*A_\Sigma
=R_0^{1/2}
\left(
I-R_0^{-1/2}A_\Sigma^*A_\Sigma R_0^{-1/2}
\right)
R_0^{1/2}.
\]

The Cut-Memory Spectral Isomorphism gives

\[
N_-(R_0-A_\Sigma^*A_\Sigma)
=N_+(B_\Sigma-I),
\]

\[
\dim\ker(R_0-A_\Sigma^*A_\Sigma)
=\dim\ker(I-B_\Sigma),
\]

and

\[
\inf_{f\ne0}
\frac{\langle f,(R_0-A_\Sigma^*A_\Sigma)f\rangle}
     {\langle f,R_0f\rangle}
=
\lambda_{\min}(I-B_\Sigma).
\]

The Fredholm determinant identity is

\[
\det_F
\left(
I-R_0^{-1/2}A_\Sigma^*A_\Sigma R_0^{-1/2}
\right)
=
\det(I-B_\Sigma).
\]

The determinant locates threshold memory. Its sign alone does not decide
positivity when \(\nu\ge2\). The lawful sign test is

\[
\lambda_{\max}(B_\Sigma)<1.
\]

\begin{center}\rule{0.5\linewidth}{0.5pt}\end{center}

\hypertarget{native-completed-weil-construction-target}{%
\section{9. Native completed-Weil construction target}\label{native-completed-weil-construction-target}}

The RH-specific theorem must be coefficientwise and cut-first.

Construct

\[
Z_{\rm Weil}f
\]

from the prime, Gamma, completion-boundary and compatibility events before any
whitening, finite Gram fit or classical explicit-formula identification.

Then prove:

\hypertarget{moment-identity}{%
\subsection{Moment identity}\label{moment-identity}}

\[
S_{0,-}^{\rm full}=Z_{\rm Weil}^*Z_{\rm Weil},
\]

\[
L_\partial(f)=\langle p_\partial,Z_{\rm Weil}f\rangle,
\]

and hence

\[
S_{0,-}^{\rm full}-L_\partial^*L_\partial
=Z_{\rm Weil}^*QZ_{\rm Weil}.
\]

\hypertarget{native-five-label-observer}{%
\subsection{Native five-label observer}\label{native-five-label-observer}}

Construct \(T_\Sigma\) directly from the five declared target labels and prove
that the adverse target-relevant bridge is faithful:

\[
QZ_{\rm Weil}
=T_\Sigma^*T_\Sigma QZ_{\rm Weil}
\]

on that module.

\hypertarget{favorable-hidden-memory}{%
\subsection{Favorable hidden memory}\label{favorable-hidden-memory}}

Prove that any target-null cut-memory complement contributes only a positive
square.

If these statements hold, the source and boundary are no longer independent,
the Cut-Square Identity is automatic, and the target-relevant spectral burden
reduces to at most five dimensions.

\begin{center}\rule{0.5\linewidth}{0.5pt}\end{center}

\hypertarget{non-circularity-rules}{%
\section{10. Non-circularity rules}\label{non-circularity-rules}}

The theorem is invalid if any of the following occurs:

\begin{verbatim}
Z_Weil is fitted from the desired completed-Weil Gram;
P or Q is chosen after seeing the target sign;
T_Sigma is obtained by diagonalizing the final operator;
rank five is inferred only from a floating singular-value threshold;
the hidden complement is omitted rather than proved favorable;
the classical explicit formula is used to define the native event lift.
\end{verbatim}

The construction is lawful only when the cut creates the objects before the
sign question is asked.

\begin{center}\rule{0.5\linewidth}{0.5pt}\end{center}

\hypertarget{claim-boundary}{%
\section{11. Claim boundary}\label{claim-boundary}}

\begin{verbatim}
cut-generated decomposition S=R+D                       PROVED
cut-generated signature F=R-D                           PROVED
rank-one first/second-moment covariance                  PROVED
exact exponential cut covariance                        PROVED
target-faithful rank-at-most-five reduction              PROVED
hidden favorable complement criterion                    PROVED
finite seam spectral/inertia/determinant reduction       PROVED

native coefficientwise Z_Weil                            NEXT
native first/second moment identification                 NEXT
native five-label observer T_Sigma                       NEXT
outward target-faithfulness and tail certificate          NEXT
RH                                                       NOT CLAIMED
\end{verbatim}

\begin{flushright}\small\texttt{RKF/theorum/23\_native\_cut\_generated\_object\_theorem.md}\end{flushright}
